%==============================================================================
% Designing Software-Intensive Enterprise Systems:
% A Critical Review of Eight Systems Traditions and a Research Agenda
% for Measuring Design Work
%
% Target venue: Systems (MDPI), Review article
%
% NOTE ON CLASS FILE. This draft uses the standard `article' class so that it
% compiles anywhere without the MDPI template installed. To convert for
% submission, download the Systems LaTeX template from
% https://www.mdpi.com/authors/latex and move the frontmatter into the
% \Title / \Author / \abstract / \keyword macros of mdpi.cls. The section
% order below already matches MDPI's expected layout, including the required
% back matter.
%==============================================================================

\documentclass[11pt,a4paper]{article}

\usepackage[T1]{fontenc}
\usepackage[utf8]{inputenc}
\usepackage[margin=1in]{geometry}
\usepackage{amsmath,amssymb}
\usepackage{graphicx}
\usepackage{booktabs}
\usepackage{longtable}
\usepackage{array}
\usepackage{url}
\usepackage{xcolor}
\usepackage[numbers,sort&compress]{natbib}
\usepackage{hyperref}
\usepackage{tikz}
\usetikzlibrary{positioning,arrows.meta,fit,backgrounds}
\usepackage{setspace}
\usepackage{lineno}

\hypersetup{colorlinks=true,linkcolor=blue!50!black,citecolor=blue!50!black,urlcolor=blue!50!black}

\newcommand{\rq}[1]{\textbf{RA#1}}

\title{\bfseries Designing Software-Intensive Enterprise Systems:\\
A Critical Review of Eight Systems Traditions and a\\
Research Agenda for Measuring Design Work}

\author{
Steven M. Morgan\thanks{Corresponding author: \url{smorgan@binghamton.edu}}\\
\small Department of Systems Science and Industrial Engineering\\
\small Binghamton University, State University of New York, Binghamton, NY 13902, USA
\and
Jeff Daniels\\
\small Graduate School of Technology\\
\small University of Maryland Global Campus, Adelphi, MD 20783, USA
}

\date{}

\begin{document}

\onehalfspacing
\linenumbers

\maketitle

%==============================================================================
\begin{abstract}
\noindent
Software-intensive enterprise systems are designed by one community and
evaluated by another. Eight distinct traditions---systems of systems (SoS),
enterprise architecture (EA), enterprise systems engineering (ESE), complex
adaptive systems (CAS), design thinking, systems thinking, systems science, and
system dynamics---each offer an account of how such systems should be conceived,
bounded, and improved, while the software engineering community independently
measures the artifacts those designs produce. This article reviews the eight
traditions as a single comparative landscape rather than as separate
literatures, and asks a question that none of them answers on its own: how would
we know, from evidence rather than from assertion, whether design work improved
a system? We find three failure modes that recur across all eight traditions.
First, a \emph{boundary blind spot}: every tradition acknowledges that emergent
behavior concentrates where constituent systems and the humans operating them
meet, and every tradition then delegates that region to some other discipline.
Second, a \emph{measurement deficit}: proposals to quantify design value depend
almost entirely on expert elicitation---stakeholder scoring, retrospective case
narration, participatory model building---which does not scale, cannot be
replicated, and is not independent of the designers being evaluated. Third,
\emph{evidentiary asymmetry}: the traditions are supported overwhelmingly by
successful case narratives, with no comparable corpus of negative results,
leaving effect sizes unknown. We argue that these are one problem, not three:
the field lacks an observational instrument for design activity. We then set out
a five-part research agenda that treats the operational record of software
projects---issue trackers, commit histories, and design discussions---as the
missing instrument, and specify what each tradition would need to commit to for
its claims to become falsifiable. The agenda is illustrated with the emerging
practice of design mining, the automated classification of design-related
discussion in engineering artifacts, and we discuss the construct-validity
conditions under which such measurements could legitimately stand in for design
work.
\end{abstract}

\vspace{0.5em}
\noindent\textbf{Keywords:} systems of systems; enterprise architecture;
enterprise systems engineering; complex adaptive systems; systems thinking;
system dynamics; design thinking; socio-technical systems; software
architecture; design mining; research agenda
\vspace{1em}

%==============================================================================
\section{Introduction}
%==============================================================================

There is a persistent gap between how information technology systems are
designed and how they are implemented. The gap does not typically present as
minor deviation. Critical requirements are misconstrued entirely, and users are
left to sustain workarounds that quietly become the operational system. This is
not a niche observation. The Standish Group's CHAOS report, drawing on roughly
5{,}000 software projects annually, found that between 2010 and 2015 only 36\%
to 41\% of software projects were judged successful, with approximately 20\%
classified as outright failures and the remainder challenged~\citep{standish2015}.
Worldwide IT spending was projected to exceed \$5 trillion in
2024~\citep{gartner2024}. A recent synthesis of three decades of project success
surveys confirms that the definitional instability of ``success'' has itself
been a barrier to cumulative knowledge~\citep{loureiro2024}, and practitioner
surveys continue to report that more than a third of information systems
projects only partially succeed or fail outright~\citep{pereira2021}. By
contrast, engineering projects with physical constraints---mining, power
generation, renewables---report markedly higher success rates on the same
on-time, on-budget, in-scope criteria.

The disciplines responsible for closing this gap have not been idle. Over five
decades, at least eight identifiable traditions have developed accounts of how
complex socio-technical systems should be conceived, decomposed, bounded, and
improved: systems of systems (SoS) architecture, enterprise architecture (EA),
enterprise systems engineering (ESE), complex adaptive systems (CAS), design
thinking, systems thinking, systems science, and system dynamics. Each has
produced frameworks, standards, professional bodies, and case evidence. Each
explicitly acknowledges that designing systems not constrained by the physical
world is harder than designing those that are, and that the difficulty is
largely human in origin.

What these traditions have not produced is a shared, external way of telling
whether their prescriptions work. Enterprise architecture, the tradition with
the most direct claim on industrial practice, offers the clearest illustration:
after decades of adoption, it remains difficult to articulate how EA creates
value~\citep{gong2019,saintlouis2017}, and the serious attempts to quantify the
value of architectural decisions have relied on stakeholder
scoring~\citep{moore2003} or qualitative workshop
consensus~\citep{nord2014}---instruments that require the continuous
participation of the very people whose work is being evaluated.

Meanwhile, an entirely separate research community has spent the same period
building quantitative instruments over the operational residue of software
work. Defect prediction, effort estimation, issue classification, and
architecture-degradation detection all operate on version control histories,
issue trackers, and developer discussion at a scale of hundreds of thousands of
records~\citep{aleti2013,tawosi2022,zangari2023,herold2020}. This community
measures prolifically but seldom asks the design-theoretic question: it predicts
that defects will occur without asking whether an architectural decision caused
them.

The two communities are addressing the same underlying phenomenon from opposite
ends and are not meeting in the middle. This article is about that gap.

\subsection{Contribution and structure}

This is a critical review with a constructive aim. We make four contributions.

\begin{enumerate}
\item \textbf{A comparative landscape.} We review eight systems design
traditions against a common analytical frame---unit of analysis, treatment of
human agency, stance toward emergence, evidence base, and measurement
instrument---rather than reviewing each in its own terms
(Sections~\ref{sec:foundations}--\ref{sec:pm}). To our knowledge these eight
literatures have not previously been placed side by side on a measurement
criterion.

\item \textbf{Three cross-cutting failure modes.} We identify and evidence a
boundary blind spot, a measurement deficit, and an evidentiary asymmetry that
appear in all eight traditions, and argue they are surface expressions of a
single missing capability (Section~\ref{sec:synthesis}).

\item \textbf{A diagnosis.} We argue the field lacks an observational instrument
for design activity---something that can detect design work in progress, at
scale, without the cooperation of the designers---and we specify what such an
instrument would have to satisfy (Section~\ref{sec:gap}).

\item \textbf{A five-part research agenda.} We set out concrete, falsifiable
research directions that would convert each tradition's central claim into a
testable proposition, and identify what each tradition must commit to in
exchange (Section~\ref{sec:agenda}).
\end{enumerate}

Section~\ref{sec:method} states the scope and method of the review and is
explicit about what kind of review this is and is not.
Section~\ref{sec:threats} discusses the limitations of the review itself, and
Section~\ref{sec:conclusions} concludes.

%==============================================================================
\section{Scope and Review Method}
\label{sec:method}
%==============================================================================

This is a \emph{critical narrative review} in the sense used in the systems
literature: an argued synthesis organized around a thesis, not an exhaustive
enumeration organized around a search protocol. We state this plainly because
the distinction bears on how the findings should be read.

\subsection{What was reviewed}

The corpus was assembled around the question \emph{how does each tradition
justify the claim that its form of design improves system outcomes?} Sources
were drawn from four pools:

\begin{itemize}
\item \textbf{Defining and standardizing sources} for each tradition---the
ISO/IEC/IEEE 21839 definition of systems of systems~\citep{iso21839}, Maier's
architecting principles~\citep{maier1998}, Holland on complex adaptive
systems~\citep{holland1992}, Klir on systems science~\citep{klir2001}, Forrester
on system dynamics~\citep{forrester2009}, the SEBoK treatment of enterprise
systems engineering~\citep{sebok_ese2024}, and the Zachman and TOGAF
frameworks~\citep{zachman_nd,bente2012}.

\item \textbf{Existing systematic reviews and surveys within each tradition},
which were used as coverage proxies rather than re-executed: Saint-Louis et
al.'s review of 174 EA papers~\citep{saintlouis2017}, Gong and Janssen's review
of EA value creation~\citep{gong2019}, Nielsen et al.'s survey of SoS modeling
techniques~\citep{nielsen2015}, Arnold and Wade's survey of systems thinking
definitions~\citep{arnold2015}, Lapalme et al.'s grand-challenges
synthesis~\citep{lapalme2016}, and Aleti et al.'s systematic review of 188
software architecture optimization papers~\citep{aleti2013}.

\item \textbf{Application and case papers} that report the traditions in use,
selected to span domains---healthcare and telecommunications~\citep{kotze2016},
defense and national security~\citep{morgan2023}, business process
management~\citep{fleischmann2016}, health promotion~\citep{vanbeurden2011}, and
commercial product organizations~\citep{dunne2018,monat2020}.

\item \textbf{Computational software-engineering work} that measures the
artifacts of design at scale, included specifically to test whether the
measurement deficit identified in the systems traditions is also present
there~\citep{aleti2013,mahadi2021,tawosi2022,zangari2023,kleebaum2021,sharma2021,dhaouadi2023}.
\end{itemize}

\subsection{What this review is not}

We did not execute a registered search protocol with reproducible strings,
database-by-database yields, or PRISMA-style inclusion counts. Readers should
therefore treat the coverage claims as \emph{argued} rather than
\emph{exhaustive}. Where we assert that a tradition does not address a question,
that assertion is grounded in the defining sources and in the existing
systematic reviews of that tradition, both of which are cited at the point of
claim; it is not grounded in a negative search result. Section~\ref{sec:threats}
returns to what this costs the argument.

\subsection{The comparative frame}

Each tradition is examined against five questions, chosen because they are the
minimum needed to compare traditions that do not share vocabulary:

\begin{enumerate}
\item \textbf{Unit of analysis.} What object does the tradition take as the
system to be designed?
\item \textbf{Human agency.} Are people modeled as components, as stakeholders,
as designers, or not at all?
\item \textbf{Stance toward emergence.} Does the tradition treat emergent
behavior as controllable, as mitigable, or as irreducible?
\item \textbf{Evidence base.} What kind of evidence supports the tradition's
central claim?
\item \textbf{Measurement instrument.} What, concretely, would one measure to
falsify the tradition's claim?
\end{enumerate}

Table~\ref{tab:landscape} in Section~\ref{sec:synthesis} presents the completed
frame.

%==============================================================================
\section{Foundations: Complex Systems, Enterprises, and Organizations}
\label{sec:foundations}
%==============================================================================

Before comparing traditions it is necessary to fix three terms that they use
differently, because much of the apparent disagreement between them is
definitional.

\textbf{Complex systems} exhibit emergent and at times chaotic properties and
resist characterization because of the number of interacting properties that
govern their behavior~\citep{siddique2024}. Critically, complex systems exist
independently of human-defined constructs; a system does not become complex
because we named it so.

\textbf{Enterprise systems} are socio-technical systems designed and created to
serve a purpose for an organization~\citep{xu2016}. An organization may be an
enterprise system or may comprise several. An enterprise system may or may not
also be a complex system, depending on whether it exhibits emergent or chaotic
behavior. This conditional relationship matters: much of the systems design
literature implicitly assumes that enterprise systems are complex systems and
imports the attendant epistemic humility, while much of the practitioner
literature implicitly assumes they are merely complicated and imports the
attendant confidence that sufficient analysis will yield a correct design.

\textbf{Enterprises} can be distinguished from \textbf{organizations}, though
the two are frequently conflated. Organizations are well-defined structures
describing relationships between individuals and groups. Enterprises may be
self-organizing, may lack a single point of control, and therefore may not be
organizations at all. The decisive difference is teleological: enterprises exist
to create value, whereas an organization need not
~\citep{sebok_ese2024,kotze2016}. A tradition that takes the organization as its
unit of analysis and a tradition that takes the enterprise as its unit will
disagree about scope even when they agree about method.

With these distinctions in place, the recurring structure of the field becomes
visible. Systems of systems, enterprise architecture, enterprise systems
engineering, systems thinking, and design thinking all attempt to capture the
intent of a system and to shape it so that it performs well. All acknowledge a
distinctly human element. All can point to successful case studies. All can also
point---though they rarely do in print---to attempts that were not productive.
The sections that follow take each in turn.

%==============================================================================
\section{Systems of Systems}
\label{sec:sos}
%==============================================================================

ISO/IEC/IEEE 21839 defines a system of systems as a ``set of systems or system
elements that interact to provide a unique capability that none of the
constituent systems can accomplish on its own,'' where a constituent system is
``an independent system that forms part of a system of systems''~\citep{iso21839}.

Maier's foundational treatment deliberately declines to build a taxonomy,
arguing that any proposed classification is rife with ambiguity and conflict and
that specific examples expose subtleties the categories cannot
hold~\citep{maier1998}. In place of taxonomy, Maier offers architecting
principles: focus on communications, policy control, interfaces, and
cooperation. The prescription that follows is consequential---in environments
requiring many distinct groups to collaborate, the appropriate move is to
encourage cooperative effort, because cooperation reduces dependency among
constituent parts.

This is in explicit contrast to building a monolith: a large, complicated system
whose parts provide marginal standalone utility yet are required for the whole
to function. Most components of an automobile are useless in isolation; a
radiator cools an engine and does little else. The internet is the canonical
positive example, built on standards for communication, interfaces, cooperation,
and policy, with Maier's principles visible at each layer and preserved as
higher layers were built on the base technical layers.

Two limitations recur. First, the modeling literature has concentrated on
engineered rather than socio-technical systems. Nielsen et al. surveyed the
modeling techniques in use and enumerated the academic and professional bodies
advancing the field, and while they acknowledge that modeling socio-technical
systems is an important problem, they explicitly leave the engineering of those
requirements to other publications~\citep{nielsen2015}. The acknowledgment
without the treatment is the pattern this review returns to repeatedly.

Second, where socio-technical constraints are engaged, the result tends to be
descriptive rather than predictive. Fleischmann et al. combine SoS with
Subject-Oriented Business Process Management (SBPM), using SoS to determine
which parts of a business process can be treated as independent systems and
invoking Maier's criteria to establish that a process fragment may be a
constituent system capable of operating and managing itself~\citep{fleischmann2016}.
SBPM treats messages within a business process as encapsulation, which maps
cleanly onto the SoS notion of an interface, and it explicitly permits variation
driven by local law, IT implementation, or market dynamics. The resulting
Pattern-Based Engineering framework does consider socio-technical systems and
their constraints. What it does not provide is any account of how the model
would improve decision making under uncertain dynamics, or how it would
translate into a simulation. Its conclusion is a statement about swappability:
the framework identifies which aspects of similar systems could be exchanged.

The interoperability consequences are vivid in the defense domain. Morgan and
Wightman describe U.S. military systems designed thirty to fifty years ago that
were never intended to communicate seamlessly across service branches or with
supply chain partners~\citep{morgan2023}. Battle management systems were built
to serve individual branches, and even within a single branch they comprise
multiple lower-level environments that do not share data. ``System integration''
in practice reduces to an operations analyst physically turning to a different
monitor. The fragmentation extends past the battlespace into the defense supply
chain, where original equipment manufacturers and their tiered suppliers operate
with divergent data models, inconsistent part numbering between engineering and
enterprise resource planning systems, and no common digital thread linking
product lifecycle to field operations. Their proposed remedy---treating the
defense enterprise as the sum of all its parts, with standardized data threads
enabling an extended Observe--Orient--Decide--Act loop spanning both campaigns
and supply chain---is an architectural claim about boundaries.

\paragraph{Assessment.} SoS is tightly coupled to systems engineering and
enterprise architecture: an enterprise architecture is in effect built from an
SoS design, and systems engineering is used to design, build, and measure the
effectiveness of the result. The tradition's characteristic gap is a lack of
emphasis on the human-centric aspects of system design, and specifically on what
happens when humans interact at the boundaries between constituent systems.
Constituent systems may be operationally independent, but mission effectiveness
requires architectural coordination precisely at the boundaries---which is the
region the tradition's own frameworks leave least specified.

%==============================================================================
\section{Enterprise Architecture}
\label{sec:ea}
%==============================================================================

Enterprise architecture describes organizational and business systems, technical
systems, their boundaries, and their interactions, and produces artifacts such
as roadmaps and models to support the design and sustainability of those
systems. It is nonetheless difficult to articulate clearly how EA creates
value~\citep{gong2019,saintlouis2017}. There are more than 77 distinct EA
frameworks, largely describing similar things in different
vocabularies~\citep{bente2012}.

\subsection{A field without a definition}

Saint-Louis et al. conducted a systematic literature review of 174 papers
seeking a consensus definition of EA and concluded that EA is better understood
as a community with definitions that vary by need than as a field with a settled
object of study~\citep{saintlouis2017}. They encourage the community to clarify
EA's intent, and note that consensus is needed. For a review concerned with
measurement this is a structural finding, not a terminological quibble: a field
without an agreed object of study cannot accumulate evidence about whether
acting on that object helps.

\subsection{Value creation and its myths}

EA is not in itself a value creator; its function is to create more effective
pathways for a business to develop value. Gong and Janssen argue that EA's value
is realized through stakeholder alignment and through risk assessment preceding
enterprise initiatives, and their systematic review spans communications,
knowledge management, transformation, and strategic and political
implications~\citep{gong2019}. They conclude by naming four myths:

\begin{itemize}
\item that EA creates value;
\item that all aspects of the enterprise can be captured;
\item that the architectural vision should not be constrained by path
dependence;
\item that complexity can be reduced.
\end{itemize}

Each of these myths reappears in the other traditions reviewed here, which
suggests they are properties of the design problem rather than pathologies of EA.

Attempts to put numbers on architectural value are instructive. Moore et al.
sought to quantify architecture effort using qualitative measures such as
stakeholder quality~\citep{moore2003}; Nord et al. reviewed several approaches
to capturing design quality and impact metrics, all of which involved
qualitative scoring~\citep{nord2014}. Both approaches require direct and
sustained involvement from individuals to capture the information consistently.
This is the measurement deficit in its purest form: the only available instrument
is the practitioner's own judgment, elicited at the practitioner's own cost.

\subsection{Schools of thought and framework structure}

Lapalme distinguishes three schools of EA thought~\citep{lapalme2012}. The
\emph{Enterprise IT} perspective emphasizes the technical stack and seeks
business alignment. The second treats the enterprise as a socio-technical
system, encompassing technology, structure, policy, and other non-technical
determinants of behavior. The third considers an extended enterprise heavily
influenced by its environment, and focuses on creating an innovative, adaptive,
learning organization in a manner closely related to Senge's organizational
learning construct~\citep{senge2018}. The schools are not merely different
emphases; they imply different units of analysis and therefore different
evidence.

The operational purpose of EA is to verify that systems developed within an
organization align with higher-level organizational aspirations. This matters
most in large organizations with many departments and business units: when
individual units are given similar resources and objectives without guidance
from a coordinating function, duplication and redundancy are inevitable, and it
is genuinely difficult to determine when redundancy is beneficial and when it is
merely duplicative spending~\citep{dunne2018}. EA attempts to guide this through
policy---steering away from duplication where it does not serve the enterprise,
while preserving autonomy where units need it. Pattern-based engineering
(Section~\ref{sec:sos}) offers related guidance for how geographically or
operationally distinct departments can interact. SoS architecture is frequently
invoked by EA efforts, but there remains a lack of clarity about how to
integrate the technical and the social
constructs~\citep{lapalme2012,nielsen2015}.

The Zachman Framework organizes EA as a matrix of interrogatives and
perspectives, on the premise that the combined interrogatives, perspectives, and
their instantiations constitute an enterprise architecture. It does consider
social aspects of systems design, and its rows are explicitly \emph{not} an
organizational hierarchy---they are different interpretations, which may be
influenced by higher- or lower-level roles~\citep{zachman_nd}. Despite the
conceptual cleanliness, Zachman and others acknowledge gaps in what the model can
support~\citep{lapalme2016}.

TOGAF is broadly recognized as complementary, extending past the point where
Zachman applies. It provides tooling and describes the enterprise architect's
responsibilities in a relatable, non-linear way, and it is cyclical---loosely
derived from the waterfall software development lifecycle, which has clear gate
phases from project inception to go-live. In contrast to Zachman, TOGAF accepts
that personas and deliverables vary by context: an engineer may not need a
process specification when a use case description would serve. Both nonetheless
designate certain artifacts as necessary for deriving an architecture.

Architectures are driven by vision and business need, which are then derived
into the technical systems that make them practical; balancing change against
governance and vision is well accepted. The difficulties lie in how rigorously
the steps are followed. Implementation problems arise when emphasis falls on
artifacts and process rather than on effective delivery~\citep{bente2012}.

\subsection{Grand challenges and paths forward}

Lapalme et al. identify three grand challenges for EA: increased complexity and
uncertainty; new realities; and evolving enterprise
architecture~\citep{lapalme2016}. These converge on the pace of change in modern
business and the expectation that an IT organization will keep step with it.
Increasing emphasis on social factors, combined with the pace of change,
produces greater uncertainty. The same authors identify frameworks better
positioned to handle uncertainty and complexity, systems thinking among them
(Section~\ref{sec:st}).

The operational environment in which a technical system runs materially
constrains its architecture; path dependence limits the utopian desired
state~\citep{gong2019}. The more complex the environment and the more integrated
the system, the harder it is for stakeholders to understand the intended
solution. Two countermeasures are proposed. \emph{System co-evolution} defines
the system as the system is developed. \emph{Increasing the involvement of
architecture in the problem space}---pushing the intended designed system and
its interfaces into problem framing rather than solution specification---is
argued to address the shortcomings of earlier techniques~\citep{cole2006}.

Rhodes, Ross, and Nightingale recognize that working in an SoS requires analysis
of many attributes affecting the creation and change of socio-technical
systems~\citep{rhodes2009}. As in other frameworks, context must be weighed
rather than any single aspect; myopic focus on one aspect produces a challenged
system. Policy, strategy, knowledge, product, external factors, and technology
must all be assessed for their impact on the whole. Their central observation is
that the ``to-be'' design is usually far from the ``as-is,'' and their remedy is
an epoch-based approach in which the system is gradually defined and
intermediate states are analyzed during development. Defining epochs as points
in time at which the system changes provides a means of soliciting feedback on
implementation progress, and this genuinely fills part of the space between
proposed utopia and current gap.

\paragraph{Assessment.} The epoch-based approach nonetheless illustrates the
tradition's limit. The factors used to define the epochs may be orthogonal to
what is actually intended; laying out intermediate steps improves implementation
but remains bounded by what human cognition can hold. The approach does not
account for emergent behavior---it acknowledges that emergence exists.

%==============================================================================
\section{Enterprise Systems Engineering}
\label{sec:ese}
%==============================================================================

Enterprise systems engineering combines traditional systems engineering with
strategic management practice to align system development and design with
desired business outcomes. Where systems engineering focuses on developing a
specific product and the supporting structures that enable it, ESE treats the
organization itself as a system and holds that the system must produce
value~\citep{kotze2016}. Producing or enhancing that value, and designing the
enterprise to do so, is ESE's primary responsibility.

The relationship between ESE and EA is contested and depends on organizational
context. Kotze and Smuts argue that ESE \emph{contains} EA, adding the shaping
of organizational vision and the definition of value outcomes~\citep{kotze2016}.
TOGAF makes structurally similar arguments on behalf of EA, with less emphasis
on strategic management practice. The substantive point beneath the demarcation
dispute is uncontroversial: if a system is being designed and implemented,
business goals and constraints must be considered for the system to be of value.

Kotze and Smuts review two cases in which information technology systems were
designed to improve operations---a South African healthcare system and a
telecommunications system---applying ESE to shape program outcomes. They claim
success and offer lessons learned from each case, and those lessons re-emphasize
the merits of ESE.

\paragraph{Assessment.} This is the evidentiary asymmetry in miniature. Two
cases, selected by the proponents, analyzed retrospectively by the proponents,
yielding lessons that confirm the approach. Nothing in the study design could
have produced a disconfirming result. We are not suggesting the cases are
misreported; we are observing that the tradition has no mechanism by which a
case could count against it, and therefore no way to estimate how large its
effect is or under what conditions it disappears.

%==============================================================================
\section{Complex Adaptive Systems}
\label{sec:cas}
%==============================================================================

Complex adaptive systems are characterized by three properties: evolution,
aggregate behavior, and anticipation~\citep{holland1992}. CAS can be understood
as systems that have no clear governing rules yet act toward common goals. The
immune system is the accessible example: it is effective at reducing threats to
the body, yet is not understood well enough to be fully recreated or repaired
unconditionally, and it adapts in anticipation of future threats. Comparable
complexity exists in supply chains and many other domains, and CAS is an attempt
to group and describe systems of this kind.

CAS has no single governing rule set and no clear representative equations.
There are equations that can simulate complex and emergent behavior---partial
differential equations, cellular automata, models of bird
migration~\citep{stevens2021}---but no generic rules for deriving such elegant
descriptions for an arbitrary CAS. Modeling and simulation has attempted to fill
the gap, and known limitations remain in what those models can do.

The tradition has been taken up in management development as a means of
understanding organizational maturation, on the premise that an organization is
not subject only to the constraints and boundaries it explicitly
adheres to~\citep{luoma2006}. Management development is conflicted by divergent
directions and implementations: what a manager seeks to accomplish may conflict
with the direction they have been given, and that conflict creates inherent
inefficiency. Luoma's Four Arenas are descriptive states developed using CAS as
a framework, and the argument is that the four states help acknowledge
complexity and the need for better techniques in management development.

CAS relates closely to the study of complexity and its interpretations---both
Holland and Luoma acknowledge this---and it heavily influenced Cynefin, a
sense-making framework for complexity~\citep{kurtz2003} that has since been
applied in domains as distant from software as health
promotion~\citep{vanbeurden2011}.

\paragraph{Assessment.} CAS is the one tradition in this review that is
epistemically honest about emergence: it treats emergent behavior as
irreducible rather than as controllable or mitigable. That honesty is purchased
at the cost of prescriptive power. CAS explains why enterprise design is hard;
it does not tell a practitioner what to do on Monday, and it offers no
measurement instrument of its own.

%==============================================================================
\section{Design Thinking}
\label{sec:dt}
%==============================================================================

Design thinking addresses some of the social constructs that systems engineering
and EA underserve, though it is arguably an implementation process rather than a
theory of systems. It begins by understanding and empathizing with a problem;
once the challenge is understood, brainstorming generates candidate solutions; a
solution is selected, prototyped, and tested~\citep{dam2024}. One can argue that
the customer-focused mindset mitigates the concerns raised
above~\citep{dunne2018}. Design thinking places particular emphasis on the
development of a solution. Its roots trace to industrial design in the 1920s,
and it has taken hold within agile development methodologies because it
encourages iterative improvement.

Its relationship to systems thinking is unsettled---the two have been
characterized as identical, distinct, integrated, and interactive. The clearest
delineation is that systems thinking emphasizes the relationships and boundaries
of a system while design thinking emphasizes cognitive
solutions~\citep{greene2017}. The origins differ and both have been applied
successfully.

Design thinking has been broadly adopted and its success rate has been
challenged. Organizations designed for efficiency do not create the space needed
to empathize toward optimal solutions; their purpose is to convert raw materials
into refined output. Design thinking is only loosely accepted as a process and
contends with cultural norms in such organizations.

The method has shown promise in think tank and laboratory settings that are not
graded strictly on output and are encouraged to innovate. This produces a
paradox. Organizations strive to be innovative, but without autonomy they do not
thrive; businesses must be efficient to survive, and efficiency is often
achieved through strict process adherence. The paradox can be overcome, but
environment and leadership must align and act promptly. Where executive
leadership creates the appropriate culture, design thinking can
thrive~\citep{dunne2018}.

\paragraph{Assessment.} Design thinking's dependency on organizational culture
is stated as a boundary condition but never operationalized. ``Leadership
created the appropriate culture'' is not a measurable antecedent, which means
that any failure of design thinking can be attributed post hoc to insufficient
cultural support. This is unfalsifiable in the strict sense, and it is a
recurring structural weakness rather than a criticism of any individual study.

%==============================================================================
\section{Systems Thinking}
\label{sec:st}
%==============================================================================

Systems thinking adopts a causal perspective on an environment and the factors
affecting it, encouraging a holistic and broad view and challenging myopic ones.
It shares much with other approaches to studying systems while sustaining its
own following. It differs from systems engineering, enterprise architecture, and
SoS architecture in that it does not emphasize aspects of system \emph{design};
it emphasizes understanding systems as they exist.

Arnold and Wade define it as ``a set of synergistic analytic skills used to
improve the capability of identifying and understanding systems, predicting
their behaviors, and devising modifications to them in order to produce desired
effects''~\citep{arnold2015}. Their survey of prevailing definitions evaluates
each against a verification test---less a rigorous test than a set of criteria
any accepted definition should satisfy---and their own definition combines the
aspects surveyed, adjusted to pass that test. The result is a widely accepted
definition reflecting the many aspects considered critical to a thorough
understanding of a system.

Senge uses \emph{The Beer Game} to introduce the problems created by not
considering systemic behavior. A convenience store sells a certain number of
cases each week and continues to order the same number. The distributor receives
orders from all the stores in its region and continues to ship and receive on
behalf of the brewery, to which it places its own orders. The brewery supplies
distributors on the basis of consistent demand. All three actors share the same
motive---to buy and sell beer---yet none is concerned with the problems the
others face elsewhere in the system. A sudden surge in demand is introduced and
participants must adjust to the new dynamics. Through the game, students learn to
see a system as the whole of its parts and the relationships among them.

Many difficulties in thinking systemically are not difficulties of capability
but of vocabulary: people do not know how to communicate in the appropriate
terms. One can easily describe one's response to being scalded by hot water; it
is much harder to articulate a problem with many contributing factors. Causal
loop analysis using reinforcing and balancing loops is a simple tool for doing
so, and outlining relationships among stocks and flows can help build
organizations genuinely capable of adopting and designing effective systems.
Senge's fuller argument holds that alongside systems thinking one must
continually develop personal mastery, adapt mental models, build shared vision,
and practice team learning, and that organizations applying these principles can
excel in challenging circumstances~\citep{senge2018}.

Monat, Amissah, and Gannon apply systems thinking to business problems across a
range of contexts, though their perspective is retrospective~\citep{monat2020}.
They provide a robust literature base for the claim that systems thinking is not
well adopted in industry and present examples illustrating failures that
followed from not adopting systems thinking methodologies, often in conjunction
with overemphasis on products. Soft Systems Methodology, Hard Systems
Methodology, and Critical Systems Methodology are reviewed to contextualize how
systems thinking is applied, yet a clear means of adopting them within an
organization remains absent. They identify common failure modes in which systems
thinking would have yielded a more complete solution, and recommend steps for
advancing systems thinking in business model application. They explicitly call
out the need to design a \emph{user experience system}: developing the support
structures surrounding a product can matter as much as the product itself, and a
product superior in isolation will not sustain success without them.

Applying systems thinking to business models differs from modeling business or
technical systems. Monat et al. support their findings with a broad reference
base and refer specifically to Sterman for depth on modeling; they include
causal loop analyses of the systems they review, a tool Sterman uses to develop
models of business dynamics.

\paragraph{Assessment.} Systems thinking provides a robust framework for
understanding problems that are not strictly technical, and it allows for
emergent behavior at system boundaries much as Maier does. But Monat et al.
argue these factors are more likely within the control of the system designer,
and that emergent properties can be mitigated through appropriate controls. This
review takes the opposing position: boundaries between socio-technical systems
are precisely where emergent behavior resists management, and the breadth of
methodologies addressing them does not clearly support the claim that such
behavior can be controlled.

%==============================================================================
\section{Systems Science}
\label{sec:ss}
%==============================================================================

Systems science comprises three primary elements: inter- and
transdisciplinarity, theories of systems, and systems approaches. Within it lies
a desire to define a system clearly~\citep{hieronymi2013}. Klir defines a system
as things and their relationships to other things---an accepted definition at
the highest level of abstraction---and network science, used to visualize and
study relationships among things within a system, is a subfield of systems
science~\citep{klir2001,hieronymi2013}.

Klir further argues that for a field to be a science there must be an agreed
object of study and a corpus of knowledge, and defines systems science as ``a
science whose domain of inquiry consists of those properties of systems and
associated problems that emanate from the general notion of
systemhood''~\citep{klir2001}. Read against Saint-Louis et al.'s finding that EA
has no consensus definition~\citep{saintlouis2017}, Klir's criterion is a
pointed one.

Systems science offers a contrast to enterprise architecture and its influencing
fields that is not as clean as it first appears. Hieronymi argues that the many
streams defining and exploring systems create
inefficiencies~\citep{hieronymi2013}. EA is focused on system design; systems
science integrates the phenomenological sciences from a mathematical
perspective, while EA takes cognitive, social, and technological systems as
inputs. A significant gap nonetheless persists between how a systems scientist
would seek to understand and define a system and how an enterprise architect
would do the same.

The influencing fields are similar for both. Yet, as Klir and Hieronymi would
argue, because the respective fields are not genuinely practicing
transdisciplinary approaches, they have diverged in method. This is the
underlying motivation of the present review: system design and software
engineering, which at their core seek to solve the same kinds of problems, have
diverged in theory and in practice, and both fields can be advanced by
examining and closing that gap.

One area that has progressed markedly is the study of dynamic and complex
systems. Methods rooted in systems science---simulation, soft computing or
computational intelligence, statistical analysis---have been advanced by
complex systems research. These concepts have yet to make a meaningful impact on
the execution of enterprise system design, which is genuinely affected by all of
the phenomenological fields systems science integrates.

%==============================================================================
\section{System Dynamics}
\label{sec:sd}
%==============================================================================

Like systems science, system dynamics emphasizes interdisciplinary study. It
views the world through a causal loop lens in which problems, actions, and
results have continuously accumulating and depleting effects. Its primary tools
are stock and flow diagrams, causal loop analysis, differential equations, and
computer simulation. Complex systems in system dynamics are understood as
fifth-order, nonlinear, dynamic systems whose relationships may be
comprehensible to humans only through simulation, and the tradition recognizes
cyclical behavior in which gain in one aspect is loss in
another~\citep{forrester2009}. Where systems science is grounded in cybernetics,
general systems theory, and complexity, system dynamics emerged from
servomechanism engineering, and it seeks to analyze the mental models reflected
in policy so that policy can be modified to support system
change~\citep{lane2000}. Analysis of mental models and their policy consequences
connects to systems science without being a primary motivation there.

Forrester proposes three database types for systems modeling: mental, written,
and numerical~\citep{forrester2009}. Mental databases contain the most knowledge
and represent the dynamic mental framing of decision making. Decision processes
may not be written down, and if they are not written they almost certainly do
not exist numerically. This ordering---most knowledge in the least accessible
form---is the sharpest statement of the measurement problem anywhere in the
literature reviewed here, and it comes from within the tradition.

\subsection{Eliciting the inaccessible}

Two families of technique attempt to reach the mental database.

\emph{Management flight simulators} study decision processes experimentally,
providing controlled inputs on time series steps with feedback loops, delays,
and stochastic behavior to simulate managerial decision making. They are, in
essence, the interactive application of system dynamics to a business, built
from stock and flow diagrams and causal loop analysis. Model data can be
collected via surveys and existing organizational sources, and validated against
the same. Organizational heuristics and decision rules can be solicited and
incorporated. Once constructed and validated, the simulator lets managers test
the effect of policy changes~\citep{papageorgiou2008}. Poor performance in these
environments is attributed to feedback delays, complex feedback, and evolving
uncertainty. Researchers have also found that more accurate mental models can be
measured using flight simulators, and that those without clear mental models
struggle to produce effective causal loop analyses~\citep{kunc2016}.

\emph{Group model building} can drive organizational change. Team participation
in developing a system dynamics model increases buy-in and can lead individuals
to alter their own mental models. Community-Based System Dynamics is one form,
addressing what is viewed as the true requirement for change: collective
conviction. A leader may develop a model leading to a policy change, but if the
affected group does not agree, system behavior remains
constant~\citep{hovmand2013}. Effectiveness can be measured in two ways. One
constructs a model with a group to evaluate which change would improve it, then
after an agreed interval asks the group to rebuild the model so that change can
be measured. The other measures the actual decisions made under the prior and
proposed mental models---if a change was intended to address, for example,
chronically excessive backlogs, backlogs can be measured once the new mental
model is adopted~\citep{kunc2016}.

System dynamics has contributed to systems thinking, systems engineering,
systems science, and design thinking, evidenced by the migration of its concepts
into those fields. Behavioral operations is a further field attempting to
incorporate psychology and social factors into our understanding of systems
behavior. Within software specifically, system dynamics has been applied to
rework uncertainty in development projects~\citep{khatun2023}, which is among
the closest existing links between this tradition and the operational data we
argue for in Section~\ref{sec:agenda}.

\paragraph{Assessment.} System dynamics is the only tradition reviewed that
possesses a genuine measurement apparatus with pre/post structure and a notion
of validation. Its limitation is throughput and independence. Flight simulators
and group model building both require participants to be convened, instrumented,
and re-convened; both measure the model of the world rather than the world; and
both are administered by the analyst whose intervention is under test. They are
laboratory instruments applied to field questions.

%==============================================================================
\section{The Role of Project Management}
\label{sec:pm}
%==============================================================================

Software projects are designed through architecture but implemented through
project management, and the distinction carries consequences for what can be
measured. A project is ``the achievement of a specific objective, which involves
a series of activities and tasks which consume resources''~\citep{munns1996}. A
project manager, by contrast, is responsible for the budget and timing of a
project, and therefore in practice for what is delivered within that timing.

Project management can succeed despite a project's failure, and vice versa. A
project may deliver the intended outcomes while running over budget and late; it
may finish on time and on budget without being of meaningful value to the
client. This is a persistent source of confusion about the role. If the project
manager is to hold time and cost within predefined bounds, they must be involved
in deciding what is actually delivered---and those three factors can conflict.
Resolving that conflict is generally considered the project manager's
responsibility. The customer's desires determine purpose and objectives; a team
creates them; the individuals on the team best understand how long each
component will take, and each carries their own motives and perceptions of what
is best for themselves and the team.

In software development the two primary methodologies are agile and waterfall.
Other forms exist but are generally classifiable within one of the two; more
than twenty recognized forms of agile project management have been
catalogued~\citep{rasnacis2017}, and a similar categorization could be performed
for waterfall.

Agile project management emphasizes self-organizing teams focused on time and
quality~\citep{sliger2011,rasnacis2017}; waterfall emphasizes time and cost.
These factors are frequently in tension---only so much quality can be delivered
for a given cost within a given time frame. Agile also emphasizes team
integration far more than waterfall implies, often prescribing daily meetings,
whereas waterfall is phase-gate focused with milestones agreed among large
stakeholder groups. Agile does not prescribe project-level checkpoints; it
encourages customer interaction through work demonstrations. Waterfall places
higher emphasis on implementing the architecture into the design, and its
phase-gate approach aligns more tightly with architectural artifacts: a formal
\emph{Approval to Proceed} may require completion and signoff of the
\emph{Concept of Operations}, and a \emph{Systems Design Review} constitutes
formal acceptance of the \emph{Systems Design Document}.

\emph{This asymmetry is directly relevant to measurement.} Waterfall generates
discrete, dated, named design artifacts; agile distributes design work across
conversations, tickets, and demonstrations without producing an equivalent
record. Any instrument that measures design work by counting artifacts will
therefore systematically under-detect design in agile projects---not because
less design occurred, but because it was recorded differently. We return to this
in Section~\ref{sec:agenda}.

With so many implementation styles, selecting a methodology is
difficult. Rasnacis and Berzisa developed a process to survey teams and select a
method based on quantitative indexes of those teams, aiming to account for
internal relationships and member motivations~\citep{rasnacis2017}. Their survey
identified relationships between team members and the number of positive choices
associated with them, with questions built on Maslow's hierarchy of needs, which
lends credibility to the psychological considerations of teamwork. By selecting
a methodology from team input and tailoring the implementation, the authors
demonstrated both buy-in and increased productivity. This is one of the few
studies in the reviewed corpus that ties a design-adjacent intervention to a
measured outcome.

Project management factors heavily into the implemented system. Although it may
not be the project manager's role to create the product, the techniques employed
can hinder or enhance the team's effectiveness~\citep{munns1996}. Most remaining
factors are what the systems architect attempts to address through design.
Neither the project manager nor the systems architect is definitively involved
in creating the end objective; both frequently aim to establish boundaries that
simplify the problem being solved. Agile has demonstrated that software projects
can be delivered more effectively under that
methodology~\citep{sandle2020,rasnacis2017}, and the most successful projects
tend to be limited in size and well defined. Much of the emphasis in SoS and EA
is on large-scale systems---not because either is inapplicable to smaller
systems, but because they were designed for large ones. Project management does
account for individual behavior, agile particularly so; SoS and EA account for
those factors differently, or, like Cynefin, function as sense-making frameworks
that help clarify what an intended solution will achieve.

Project and program management are sometimes used interchangeably. Program
management is generally a scaling up: a program comprises several projects
related through strategic objectives. Projects are thought of as having clear
outcomes; programs are higher level and organizationally focused. Portfolios are
generally managed as programs with common strategic value. A program manager
delivers value realization through collective opportunities across projects; a
project manager delivers value to the project's definition. The two are
supported by different research bases, which is reasonable given their differing
emphases~\citep{weaver2010}. Program management research draws on organization
theory, product development, work design and change, and quality and
manufacturing; project management research emphasizes innovation, performance,
product development, and management~\citep{artto2009}.

In agile methodologies, project management responsibilities are shared within
the team, primarily between the scrum master and product owner. The scrum master
convenes regular discussions, resolves blockers, and maintains the process. The
product owner represents the customer, refines the backlog to establish
priority, and defines expectations for a given capability. Depending on scale,
an individual responsible for overall project management may still
exist~\citep{sliger2011}.

Depending on organization or project size there may also be a chief engineer,
chief architect, or lead engineer serving as the authority on technical design.
The chief engineer is expected to exercise expert judgment in balancing
technical implementation against project risk, often plays a significant role in
developing roadmaps from the design, and typically validates that the technical
implementation meets the design and that variations are reconciled. This role
would benefit immensely from empirical understanding of how its decisions affect
project success---and at present has essentially no way to obtain it.

%==============================================================================
\section{Synthesis: Three Failure Modes Across Eight Traditions}
\label{sec:synthesis}
%==============================================================================

Table~\ref{tab:landscape} places the eight traditions against the comparative
frame introduced in Section~\ref{sec:method}. Three patterns cut across the rows.

\begin{table}[htbp]
\centering
\footnotesize
\caption{The eight traditions against a common comparative frame. ``Falsifiable
claim?'' asks whether the tradition, as currently formulated, specifies a
measurement that could disconfirm its central proposition.}
\label{tab:landscape}
\renewcommand{\arraystretch}{1.25}
\begin{tabular}{@{}p{1.9cm}p{2.2cm}p{2.3cm}p{2.1cm}p{2.5cm}p{1.6cm}@{}}
\toprule
\textbf{Tradition} & \textbf{Unit of analysis} & \textbf{Human agency} &
\textbf{Emergence} & \textbf{Evidence base} & \textbf{Falsifiable claim?} \\
\midrule
Systems of systems &
Constituent systems and interfaces &
Operators outside the model &
Acknowledged; delegated &
Standards, principles, exemplar architectures &
No \\

Enterprise architecture &
Organization and its technical estate &
Stakeholders to be aligned &
Acknowledged; assumed reducible &
Frameworks (77+), SLRs, practitioner reports &
No \\

Enterprise systems engineering &
The enterprise as value producer &
Strategic decision makers &
Largely unaddressed &
Proponent-authored case studies &
No \\

Complex adaptive systems &
Aggregate adaptive behavior &
Agents within the system &
Irreducible &
Analogy, simulation, descriptive states &
No \\

Design thinking &
The solution and its users &
Central; empathy-driven &
Not addressed &
Practice reports; culture-conditioned outcomes &
No \\

Systems thinking &
Relationships and boundaries &
Mental models of participants &
Claimed mitigable via controls &
Retrospective cases; definitional surveys &
Weakly \\

Systems science &
Systemhood as a property &
Cognitive/social input domains &
Object of formal study &
Formal theory; mathematical integration &
Weakly \\

System dynamics &
Stocks, flows, feedback, policy &
Mental models as primary data &
Endogenous to the model &
Simulation; elicitation experiments &
Yes, in lab \\
\bottomrule
\end{tabular}
\end{table}

\subsection{Failure mode 1: the boundary blind spot}

Every tradition reviewed identifies the interface between constituent systems,
and between systems and the humans operating them, as where difficulty
concentrates. Maier's architecting principles are entirely about boundaries:
communications, policy control, interfaces, cooperation~\citep{maier1998}.
Morgan and Wightman's defense analysis locates mission failure precisely at
boundaries between branch systems and supply chain
partners~\citep{morgan2023}. EA invokes SoS to handle boundaries but cannot
specify how technical and social constructs integrate across
them~\citep{lapalme2012,nielsen2015}. Systems thinking foregrounds boundaries as
its defining concern~\citep{greene2017}.

And every tradition then hands the boundary region to someone else. Nielsen et
al. name socio-technical modeling as important and defer it to other
publications~\citep{nielsen2015}. Fleischmann et al. build a framework that
accommodates socio-technical constraints and conclude with a statement about
swappability rather than about behavior~\citep{fleischmann2016}. Rhodes et al.
acknowledge emergent behavior and proceed to a method that does not account for
it~\citep{rhodes2009}. Gong and Janssen list ``all aspects of the enterprise can
be captured'' as a myth~\citep{gong2019} without proposing what to do once it is
abandoned.

This is not carelessness. It reflects that the boundary region is where human
behavior dominates and where the traditions' instruments---models, frameworks,
reference architectures---have the least purchase. The blind spot is
instrument-shaped.

\subsection{Failure mode 2: the measurement deficit}

Every instrument for assessing design value found in this review requires the
continued, active participation of human experts. Moore et al. use stakeholder
quality scoring~\citep{moore2003}. Nord et al. survey design quality metrics and
find them uniformly qualitative~\citep{nord2014}. Kotze and Smuts narrate
outcomes retrospectively~\citep{kotze2016}. Monat et al. reconstruct failures
after the fact~\citep{monat2020}. Group model building requires convening the
affected group twice~\citep{hovmand2013,kunc2016}. Flight simulators require
constructing and validating a bespoke simulation per
organization~\citep{papageorgiou2008}.

Three consequences follow. First, \textbf{these instruments do not scale}: the
cost of measurement rises linearly with the number of systems studied, which is
why the evidence base consists of individual cases rather than populations.
Second, \textbf{they are not replicable}: a second team eliciting from the same
organization would produce different scores, and no study in the reviewed corpus
reports inter-rater reliability for an architectural value judgment. Third, and
most seriously, \textbf{they are not independent}: the people scoring the
architecture are, in the main, the people who produced it or commissioned it.
Forrester's own ordering---most knowledge in the mental database, least in the
numerical~\citep{forrester2009}---explains why the field settled here, but it
does not make the resulting evidence stronger than it is.

\subsection{Failure mode 3: evidentiary asymmetry}

The corpus is dominated by accounts of success. ESE is supported by two cases
the authors describe as successful~\citep{kotze2016}. Design thinking is
supported by settings in which it flourished, with failures attributed to
cultural conditions~\citep{dunne2018}. Systems thinking is supported by
retrospective reconstructions of failures that systems thinking would
purportedly have prevented~\citep{monat2020}---a structure that cannot
disconfirm, since the counterfactual is never run. Rasnacis and Berzisa's
methodology-selection study~\citep{rasnacis2017} is among the few that reports a
measured outcome against a pre-stated intervention.

The absence of negative results has a specific cost: \emph{effect sizes are
unknown}. Nothing in the reviewed literature supports a statement of the form
``architectural investment of this magnitude, at this point in a project, is
associated with this change in outcome.'' Practitioners are asked to invest in
design on the strength of consensus and case narrative. That may well be
correct. It is not established.

\subsection{One problem, not three}

These three failure modes have a common cause. A field that could observe design
activity independently of the designers---at population scale, repeatably, and
without requiring cooperation---would be able to study boundary behavior where
it actually occurs, would not depend on elicitation, and would generate
disconfirming as well as confirming cases as a matter of course. The three
failure modes are what the absence of such an instrument looks like from three
different angles.

%==============================================================================
\section{The Measurement Gap and What Would Close It}
\label{sec:gap}
%==============================================================================

We state the diagnosis directly. \textbf{The systems design traditions lack an
observational instrument for design work.} They possess rich theories of what
design should be and prescriptive frameworks for how to conduct it, and they
have no way to detect that it happened.

Figure~\ref{fig:gap} summarizes the position. Two communities operate on the
same phenomenon from opposite ends, and the agenda of
Section~\ref{sec:agenda} is an attempt to connect them.

\begin{figure}[htbp]
\centering
\begin{tikzpicture}[
  font=\small,
  box/.style={draw, rounded corners=2pt, align=center, inner sep=5pt, minimum height=1.5cm},
  trad/.style={box, fill=blue!6, draw=blue!45, text width=4.1cm},
  arti/.style={box, fill=orange!8, draw=orange!55, text width=4.1cm},
  gapbox/.style={box, fill=gray!8, draw=gray!55, dashed, text width=3.3cm},
  ra/.style={draw=gray!60, -{Latex[length=2mm]}, thick}
]

\node[trad] (trad) at (0,0)
  {\textbf{Systems design traditions}\\[2pt]
   \footnotesize SoS $\cdot$ EA $\cdot$ ESE $\cdot$ CAS\\
   \footnotesize Design thinking $\cdot$ Systems thinking\\
   \footnotesize Systems science $\cdot$ System dynamics};

\node[arti] (arti) at (10,0)
  {\textbf{Operational software record}\\[2pt]
   \footnotesize Issue trackers $\cdot$ commit histories\\
   \footnotesize design discussion $\cdot$ defect records\\
   \footnotesize $10^5$--$10^6$ timestamped artifacts};

\node[gapbox] (gap) at (5,0)
  {\textbf{The measurement}\\\textbf{gap}\\[3pt]
   \footnotesize No instrument links\\ \footnotesize construct to artifact};

\node[align=center, text width=4.1cm, font=\footnotesize\itshape] at (0,-1.7)
  {Rich constructs.\\ No independent, scalable measurement.};
\node[align=center, text width=4.1cm, font=\footnotesize\itshape] at (10,-1.7)
  {Abundant measurement.\\ No design-theoretic constructs.};

\draw[ra] (trad.east) -- node[above, font=\footnotesize] {what to measure} (gap.west);
\draw[ra] (arti.west) -- node[above, font=\footnotesize] {what is observable} (gap.east);

\node[align=center, font=\footnotesize, text width=11cm] at (5,-2.9)
  {\textbf{RA1} construct validity \quad $\cdot$ \quad \textbf{RA2} outcome effects
   \quad $\cdot$ \quad \textbf{RA3} temporal distribution\\
   \textbf{RA4} boundary instrumentation \quad $\cdot$ \quad
   \textbf{RA5} elicited vs.\ observed reconciliation};

\draw[ra, dashed] (gap.south) -- (5,-2.5);

\end{tikzpicture}
\caption{The measurement gap. The systems design traditions supply constructs
for what design work is and why it should matter; the operational record of
software projects supplies observations at a scale and independence the
traditions' elicitation-based instruments cannot reach. Neither side currently
provides the mapping between them. The five research directions of
Section~\ref{sec:agenda} are staged attempts to construct that mapping, with
RA1 (construct validity) as the precondition for the remaining four.}
\label{fig:gap}
\end{figure}

An adequate instrument would need four properties.

\begin{description}
\item[Independence.] It must observe design activity without requiring the
designers to report on it. Any instrument mediated by self-report inherits the
self-affirmation bias that the traditions' case literature already exhibits.

\item[Scale.] It must operate over populations of projects, not exemplars.
Effect sizes and boundary conditions cannot be estimated from case studies
selected by proponents.

\item[Temporal resolution.] It must locate design work \emph{in time}, not
merely record that it occurred. The central disputes between agile and waterfall
(Section~\ref{sec:pm}), and the epoch-based proposals of Rhodes et
al.~\citep{rhodes2009}, are claims about \emph{when} design should happen. They
cannot be tested by an instrument that only counts.

\item[Construct validity.] It must measure design work rather than a correlate
of it, and the burden of demonstrating this falls on the instrument's
proponents, not its critics.
\end{description}

Such an instrument is closer to hand than the systems literature appears to
recognize, because the parallel software engineering community has been building
its raw materials for two decades. Software projects leave a dense operational
record: issue trackers containing hundreds of thousands of tickets, version
control histories, and threaded design discussion. The TAWOS dataset alone
comprises over 458{,}000 issues and 1.6 million comments from 39 open source
projects~\citep{tawosi2022}. This record satisfies independence and scale by
construction---it is generated as a byproduct of work rather than elicited---and
it is inherently timestamped, satisfying temporal resolution.

What it does not satisfy automatically is construct validity, and that is the
hard part. A ticket is not a design decision. The gap between the two is the
research problem, and it is the subject of the next section.

The relevant computational capability already exists in adjacent form.
\emph{Design mining}---the automated classification of design-related
discussion in software engineering artifacts---has been established on Stack
Overflow data, along with the crucial finding that conclusions drawn on one
corpus do not necessarily hold on another~\citep{mahadi2021}. Rationale
extraction has been pursued in issue trackers and version control
systems~\citep{kleebaum2021}, in project email archives~\citep{sharma2021}, and
in commit messages via knowledge graphs~\citep{dhaouadi2023}. Machine learning
has been applied to architecture degradation~\citep{herold2020} and to
multi-level ticket classification~\citep{zangari2023}. Transformer
architectures~\citep{devlin2019} have made the text classification component
tractable at the scale required.

Aleti et al.'s systematic review of 188 software architecture optimization
papers marks the boundary of what has been attempted~\citep{aleti2013}. Their
review covers optimization against quality criteria---performance, security,
reliability---and constraints including cost and quality of service. Impacts on
how the \emph{agents} of the system were affected by decisions such as
prioritization were deliberately excluded. The distribution is telling: 51\% of
the optimization work targeted embedded systems and only 26\% included
enterprise systems, with metaheuristics and genetic algorithms dominating
because they suit multi-objective search. The work is valuable, and excluding
papers concerned with the interaction of the people who build and use the system
is exactly the exclusion this review is concerned with. Optimizing the build of
a system and optimizing for the interaction of its users are highly interrelated
and should not be treated as discrete.

Soft computing was defined by Zadeh as an attempt ``to mimic the ability of the
human mind to effectively employ modes of reasoning that are approximate rather
than exact''~\citep{zadeh1994}. That is a reasonable description of what an
instrument for design work must do. The techniques exist; they have not been
pointed at this question.

%==============================================================================
\section{A Research Agenda}
\label{sec:agenda}
%==============================================================================

We propose five research directions. Each is stated as a falsifiable proposition
together with the commitment the relevant tradition must make for the test to be
meaningful.

\subsection{RA1: Establish construct validity for artifact-based design detection}

\textbf{Proposition.} Design-related activity can be detected in software
engineering artifacts with agreement to expert human judgment sufficient to
serve as a measurement instrument, and that agreement is stable across projects.

\textbf{Why it is first.} Every subsequent direction depends on it. The known
obstacle is conclusion instability: classifiers trained on one corpus do not
necessarily transfer to another~\citep{mahadi2021}. Cross-project evaluation, not
within-corpus accuracy, is the relevant standard, and the appropriate protocol
is leave-one-project-out validation in which the evaluation project contributes
nothing to training.

\textbf{What must be committed.} Design mining proponents must report per-project
performance rather than pooled averages, must report agreement with multiple
independent annotators rather than a single labeler, and must characterize what
the classifier attends to. Attribution analysis that reveals a model keying on
project-specific vocabulary---container tooling, deployment platforms, plug-in
names---is evidence \emph{against} construct validity even when aggregate
accuracy is high, and should be reported as such.

\subsection{RA2: Estimate the effect of design activity on downstream outcomes}

\textbf{Proposition.} Variation in detected design activity within a project is
associated with subsequent variation in defect rates, and the association
survives controls for project size, team size, and activity volume.

\textbf{Why it matters.} This is the claim the entire field rests on and has
never tested at population scale. EA asserts value through stakeholder
alignment~\citep{gong2019}; ESE asserts it through case
narrative~\citep{kotze2016}; neither states an effect size. Defect rate is not
the only outcome worth modeling, but it is observable in the same operational
record as the design signal, which makes it the tractable starting point. The
substantial defect-prediction literature provides both baselines and covariates.

\textbf{What must be committed.} Proponents of architectural investment must
name, in advance, the direction and rough magnitude of effect they expect. A
null result must be reportable. Analyses must be pre-specified with respect to
covariates, because the number of available controls in a repository dataset is
large enough to support post hoc model selection toward a desired conclusion.

\subsection{RA3: Test the temporal claims that separate agile from waterfall}

\textbf{Proposition.} The temporal distribution of design activity across a
project's lifecycle predicts outcome variation independently of the total volume
of design activity.

\textbf{Why it matters.} This is where the systems traditions make their
sharpest empirical commitments and have never been adjudicated. Waterfall's
phase-gate structure asserts that front-loaded design, formalized in a Systems
Design Review, reduces downstream defects. Agile asserts that detailed upfront
design decays in dynamic environments. Rhodes et al.'s epoch-based
approach~\citep{rhodes2009} and Cole's system co-evolution~\citep{cole2006} are
intermediate positions. All are testable given a timestamped design signal, and
system dynamics offers the natural modeling apparatus for the feedback
structure~\citep{forrester2009,khatun2023}.

\textbf{What must be committed.} The measurement asymmetry noted in
Section~\ref{sec:pm} must be addressed head-on: waterfall projects produce
discrete design artifacts and agile projects do not, so any comparison must
demonstrate that the instrument is not simply detecting documentation practice.
This requires validating the detector separately within each methodology before
comparing across them.

\subsection{RA4: Instrument the boundary region directly}

\textbf{Proposition.} Design discussion that spans organizational or system
boundaries exhibits different characteristics---in volume, duration, resolution
rate, or downstream defect association---than design discussion internal to a
single constituent system.

\textbf{Why it matters.} This addresses failure mode 1 directly. Every tradition
locates difficulty at boundaries and none observes boundary behavior
empirically. Cross-component issue links, cross-team assignment, and
cross-repository dependencies are all recoverable from operational data, which
makes Maier's principles~\citep{maier1998} and Morgan and Wightman's defense
observations~\citep{morgan2023} testable rather than merely plausible.

\textbf{What must be committed.} SoS and EA proponents must specify what
observable consequence their boundary prescriptions are supposed to have. A
principle that predicts nothing measurable about boundary-spanning work cannot
be evaluated, and should be reclassified as heuristic guidance rather than
architectural theory.

\subsection{RA5: Reconcile elicited and observed measures}

\textbf{Proposition.} Expert judgments of architectural significance, elicited by
the established qualitative instruments, converge with artifact-derived measures
of design activity on the same projects.

\textbf{Why it matters.} This is the direction that protects against the new
instrument simply being wrong. Stakeholder scoring~\citep{moore2003},
qualitative design quality metrics~\citep{nord2014}, and group model
building~\citep{hovmand2013} encode genuine expert knowledge that no classifier
possesses. If elicited and observed measures diverge, that divergence is itself
the finding, and locating it---systematically by project type, methodology, or
domain---would advance both communities. Forrester's mental, written, and
numerical databases~\citep{forrester2009} provide the framing: artifact-based
measurement is an attempt to recover the mental database through the written
one, and RA5 asks how much is lost in transit.

\textbf{What must be committed.} The qualitative traditions must permit their
instruments to be scored for reliability. A stakeholder scoring method whose
inter-rater agreement is not reported cannot serve as a criterion measure for
anything.

\subsection{What the agenda does not claim}

We are not proposing that computational measurement replaces the design
traditions. A classifier does not know why an architectural decision was made,
cannot evaluate whether it was wise, and has no access to the strategic context
that ESE and EA are organized around. The claim is narrower and, we think,
harder to dispute: a field that cannot observe its central phenomenon cannot
accumulate knowledge about it, and the operational record of software work is
the most promising available observation point. The traditions supply the
constructs; the artifacts supply the measurements; neither is sufficient alone.

%==============================================================================
\section{Threats to the Validity of This Review}
\label{sec:threats}
%==============================================================================

\textbf{Selection.} The review is narrative rather than systematic
(Section~\ref{sec:method}). Sources were selected to illuminate a measurement
thesis, and a different frame would surface different sources. In particular, a
reader who believes the traditions' value is primarily generative---that
frameworks are useful for structuring thought regardless of whether their
effects can be measured---will find the selection unsympathetic. We think that
position is defensible and that it should be stated as a position rather than
assumed.

\textbf{Negative claims from positive evidence.} Where we assert that a
tradition does not address a question, the support is the defining sources and
the tradition's own systematic reviews, not a systematic absence search. These
assertions are therefore rebuttable by counterexample, and we would regard
well-aimed counterexamples as the most useful response this article could
provoke.

\textbf{Coverage boundaries.} Several adjacent literatures are represented
thinly or not at all: sociotechnical systems theory in the STS tradition,
organizational information processing theory, real options approaches to
architectural investment, and the technical debt literature, which addresses a
closely related measurement problem from within software engineering. Each could
strengthen or complicate the argument.

\textbf{Domain generality.} The agenda in Section~\ref{sec:agenda} is specified
for software-intensive systems, because those systems produce the operational
record the agenda depends on. Whether the same approach extends to enterprise
systems without a dense digital work record---and much of the enterprise
architecture literature concerns exactly those---is open. We would expect the
instrument to degrade gracefully rather than transfer cleanly.

\textbf{Author position.} The authors are actively pursuing the computational
direction advocated in Section~\ref{sec:agenda}. Readers should weigh the
diagnosis accordingly. We have tried to state the construct validity problem
(RA1) at full strength for exactly this reason: it is the point on which our own
program is most vulnerable.

%==============================================================================
\section{Conclusions}
\label{sec:conclusions}
%==============================================================================

Eight traditions have developed accounts of how software-intensive enterprise
systems should be designed. They differ in vocabulary, unit of analysis, and
intellectual lineage, and reviewing them side by side on a common frame shows
they converge on the same three limitations. They locate difficulty at
socio-technical boundaries and then delegate those boundaries elsewhere. They
depend on measurement instruments that require expert elicitation and therefore
do not scale, do not replicate, and are not independent of the designers under
evaluation. They are supported by a literature of successes with no comparable
record of failures, leaving effect sizes unknown after five decades.

We have argued that these are three views of one absence: there is no
observational instrument for design activity. We have specified four properties
such an instrument would need---independence, scale, temporal resolution, and
construct validity---and argued that the operational record of software
projects satisfies the first three by construction while making the fourth the
central research problem.

The five-part agenda in Section~\ref{sec:agenda} sets out what would have to be
demonstrated, and, equally important, what each tradition would have to commit
to in exchange. Those commitments are not trivial. They require proponents to
state expected effect sizes in advance, to report per-project rather than pooled
results, to permit null findings, to score their qualitative instruments for
reliability, and to specify observable consequences for prescriptions that
currently have none. A tradition unwilling to make such commitments is not
thereby refuted, but it should be understood as offering heuristic guidance
rather than empirical knowledge, and it should say so.

The divergence between systems design and software engineering analytics is not
inevitable. Both communities are addressing the same phenomenon: how deliberate
design effort changes the behavior of a complex socio-technical system. One has
the constructs and no measurements; the other has the measurements and no
constructs. Closing that gap is the contribution this article is arguing for.

%==============================================================================
% MDPI-required back matter
%==============================================================================

\section*{Abbreviations}
\begin{tabular}{@{}ll@{}}
CAS   & Complex Adaptive Systems \\
EA    & Enterprise Architecture \\
ESE   & Enterprise Systems Engineering \\
IT    & Information Technology \\
OODA  & Observe--Orient--Decide--Act \\
SBPM  & Subject-Oriented Business Process Management \\
SD    & System Dynamics \\
SoS   & System of Systems \\
TOGAF & The Open Group Architecture Framework \\
\end{tabular}

\vspace{1em}
\noindent\textbf{Author Contributions:} Conceptualization, S.M.M.; methodology,
S.M.M.; investigation, S.M.M.; writing---original draft preparation, S.M.M.;
writing---review and editing, S.M.M. and J.D.; supervision, J.D. All authors
have read and agreed to the published version of the manuscript.

\noindent\textbf{Funding:} This research received no external funding.

\noindent\textbf{Data Availability Statement:} No new data were created or
analyzed in this study. Data sharing is not applicable to this article.

\noindent\textbf{Conflicts of Interest:} The authors declare no conflict of
interest.

%==============================================================================
\bibliographystyle{unsrtnat}
\bibliography{chapter2_refs}
%==============================================================================

\end{document}
